\begin{tikzpicture}[
  cid figure,
  state/.style={draw=figline, line width=0.65pt, rounded corners=3pt, minimum height=0.82cm, align=center, inner xsep=8pt, inner ysep=4pt},
  flow/.style={-{Latex[length=1.8mm]}, thick},
  loop/.style={-{Latex[length=1.8mm]}, thick, dashed},
  note/.style={font=\sffamily\footnotesize, align=center}
]
  % Position each state from the previous node's actual border so all four
  % horizontal arrows have the same visible length even when text widens a box.
  \node[state, cid thought, minimum width=2.25cm] (need) at (0,0) {active\\information need};
  \node[state, cid runtime, minimum width=2.35cm, right=0.75cm of need] (bind) {persistent binding\\$b_j$};
  \node[state, cid fact, minimum width=2.5cm, right=0.75cm of bind] (cache) {cached value\\version + provenance};
  \node[state, cid thought, minimum width=2.45cm, right=0.75cm of cache] (project) {contextual projection\\$P_j^s$};
  \node[state, cid display, minimum width=2.35cm, right=0.75cm of project] (state) {evolving\\thought + display};

  \draw[flow] (need) -- (bind);
  \draw[flow] (bind) -- (cache);
  \draw[flow] (cache) -- (project);
  \draw[flow] (project) -- (state);

  \node[state, minimum width=2.05cm, below=1.18cm of bind] (retire) {retire binding};
  \draw[flow] (bind.south) -- node[note, right] {need inactive} (retire.north);

  \node[state, cid fact, minimum width=2.6cm, below=1.18cm of cache] (source) {external read\\or refresh};
  \draw[flow] (source.north) -- (cache.south);
  \node[note] at ($(source.south)+(0,-0.42)$) {first fetch or source-version change};

  \draw[loop] (state.north) -- ++(0,1.0) -| (project.north);
  \node[note] at ($(project.north)!0.5!(state.north)+(0,1.38)$) {need remains active:\\re-project cached value};
\end{tikzpicture}
